%======================================================================
\section{Anatomy of a Multifrequenz signal trial}
\label{sec:signal-anatomy}

Before describing how the recordings are denoised, this section
investigates what they actually contain. Each trial is a
\SI{9}{\milli\second} waveform in which acoustic, electronic, and
mechanical effects overlap in time. Identifying which samples belong
to which physical process is a prerequisite for the design decisions
that follow: where to apply blanking, which region can be used for
noise variance estimation, and what the matched filter and the
wavelet decomposition will see.

The analysis is performed on a representative cell of the dataset:
the \texttt{Box} target at \SI{75}{cm} and $0\degree$, recorded on
\texttt{Mic2}, averaged across all 100 trials in the file. The
\texttt{Mic2} channel is chosen because it is the central
microphone of the matched pair and because the \SI{75}{cm} distance
is intermediate within the protocol grid, allowing the full
sequence of signal regions to be observed within a single trial.
A reference (target-absent) recording, also averaged over 100
trials, is used throughout as a baseline to distinguish acoustic
content from hardware-driven content.

\subsection{The transmission burst}
\label{sec:transmission-burst}

\begin{figure}[tb]
    \centering
    \includegraphics[width=0.95\linewidth]{Pictures/fig_anatomy_full.png}
    \caption{Single trial of the Multifrequenz dataset
        (\texttt{Box}, \SI{75}{cm}, $0\degree$, \texttt{Mic2},
        trial 0). Shading indicates the four signal regions
        identified empirically: (A) transmission burst, samples
        0--600; (B) transducer and microphone ringdown, samples
        600--2000; (C) stationary multipath from room reflections,
        samples 2000--10000; (D) noise floor, samples 10000--18000.
        The dashed blue line at sample 2400 marks the front-end
        blanking boundary used downstream.}
    \label{fig:anatomy-full}
\end{figure}

Figure~\ref{fig:anatomy-full}
shows a single trial from this cell with four signal regions
annotated. The first \SI{0.3}{\milli\second} (samples 0 to
approximately 600) contains the transmission burst itself. The
recorded amplitude in this region reaches the limits of the
\num{12}-bit \gls{adc} dynamic range, with the analog stage close
to saturation. The burst is not pure direct-coupling
electromagnetic crosstalk: the receiver records mechanical
displacement of the microphone diaphragm in response to the
acoustic field, and the transmitter--microphone separation of only
\SI{5}{mm} on the PCB means the first wavefront arrives at the
microphone within approximately \SI{15}{\micro\second} of the
transmitter activation. The burst is therefore predominantly an
acoustic event recorded at extremely short range, on top of an
electromagnetic component that is unavoidable given the analog
proximity of the transmitter drive lines.

A spectral analysis of this region reveals the defining property of
the Multifrequenz dataset. The three transmitters fire
simultaneously, each at its assigned carrier frequency of
\SI{40}{kHz}, \SI{50}{kHz}, or \SI{60}{kHz}. The recorded burst at
the microphone therefore contains all three carriers concurrently.
After a Hann-windowed Fourier transform of the burst region, the
three carrier peaks are within \SI{3.5}{dB} of each other in
magnitude, with the \SI{40}{kHz} component the strongest in this
dataset --- consistent with the receiver's larger response at
frequencies near its \SI{39}{kHz} resonance.
This simultaneous excitation determines several downstream design
choices: the matched filter must process each carrier band
separately (Chapter~\ref{chapter:tof}), and the wavelet
decomposition in Section~\ref{sec:wavelet-denoising} must either
preserve all three bands or address them separately depending on
the branch used.

The envelope of the burst shows a max-to-mean ratio of
approximately \num{1.7}, which is intermediate between a perfectly
rectangular window (ratio \num{1.0}) and a Hann window (ratio
\num{2.0}). The observed shape is a product of the drive waveform
and the impulse response of the transducer; the exact window shape
of the drive signal cannot be inferred unambiguously from the
recording alone, and is not necessary for the analysis that follows.

\subsection{Transducer ringdown}
\label{sec:ringdown}

After the transmitter drive cuts off at approximately
\SI{0.3}{\milli\second}, the recorded amplitude decays exponentially
as the piezoelectric transmitter element and the microphone
diaphragm continue to oscillate at their respective resonant
frequencies. For a resonator with quality factor $Q$ at frequency
$f_0$, the envelope decays as $e^{-t/\tau}$ with time constant
$\tau = Q / (\pi f_0)$. For the Knowles SPH18C3LM4H-1 microphone,
Table~2 of the datasheet \citep{Knowles2020} specifies a resonant
peak at \SI{39}{kHz} but does not quote a quality factor. A
quality factor can be estimated from the curve in Figure~7 of the
datasheet by reading the $-3$~dB points around the resonance peak;
this estimate is approximately $Q \approx 2.5$, although the
limited resolution of the figure makes it uncertain by a factor of
roughly two. Adopting $Q \approx 2.5$ as a working value gives
\begin{equation}
    \tau \approx \frac{2.5}{\pi \cdot \SI{39}{kHz}}
            \approx \SI{20}{\micro\second}.
    \label{eq:tau-ringdown}
\end{equation}
The hardware engineer's own expectation, given that the
\SI{40}{kHz} carrier sits at the resonance peak, was that the
microphone should not exhibit substantial ringing in operation.
At the time constant in
Equation~\eqref{eq:tau-ringdown}, the ringdown amplitude falls by
a factor of $e^{-50}$ in just \SI{1}{\milli\second}, so a pure
exponential ringdown from the burst peak should reach the noise
floor in well under \SI{1}{\milli\second}, in line with that
expectation.

\begin{figure}[tb]
    \centering
    \includegraphics[width=0.95\linewidth]{Pictures/fig_anatomy_decay.png}
    \caption{Reference envelope decay (Hilbert magnitude, averaged
        over 100 trials, blue) compared to a pure exponential
        decay with $Q \approx 2.5$ at $f_0 = \SI{39}{kHz}$
        ($\tau \approx \SI{20}{\micro\second}$, red dashed),
        anchored at the end of the transmission burst at sample
        600. The empirical envelope tracks the theoretical decay
        only for the first $\sim$\SI{0.5}{\milli\second}; beyond
        that, the empirical curve lies an order of magnitude or
        more above the theoretical line, and from approximately
        \SI{1}{\milli\second} onwards it ceases to be monotonic.
        The non-monotonic behaviour is incompatible with a single
        ringdown process and identifies the
        \SIrange{1}{5}{\milli\second} region as stationary
        multipath rather than continued ringdown.}
    \label{fig:anatomy-decay}
\end{figure}

Figure~\ref{fig:anatomy-decay}
compares the empirical envelope of the reference recording (Hilbert
magnitude, averaged over 100 trials) against this theoretical
exponential decay, anchored at the end of the transmission burst.
The empirical envelope tracks the pure-ringdown prediction closely
for approximately the first \SI{0.5}{\milli\second}; thereafter, the
empirical envelope lies an order of magnitude or more above the
theoretical line. From approximately \SI{1}{\milli\second} onwards
the empirical envelope ceases to decay monotonically, exhibiting
local maxima and minima rather than a clean exponential tail. This
non-monotonic behaviour is incompatible with a single ringdown
process and motivates a separate identification of what dominates
the recording between \SI{1}{\milli\second} and the noise floor.

\subsection{Stationary multipath from room reflections}
\label{sec:multipath}

Between samples 2000 and 10000 (approximately
\SI{1}{\milli\second} to \SI{5}{\milli\second}), the envelope shows
bumps that rise and fall on time scales of a few hundred
microseconds, with overall amplitude decreasing slowly toward the
noise floor. These bumps are present in the reference recording,
where no target is placed in front of the array, and they are
essentially identical across the 100 trials of any one file --- the
across-trial standard deviation of the envelope in this region is
small compared to its trial-averaged mean.

Two observations identify this content as stationary multipath
from the test setup rather than acoustic returns from the target
or from continuing transducer resonance. First, it persists
identically in the absence of a target, ruling out target-related
echoes. Second, its envelope is not monotonic, ruling out a single
ringing resonator: the local maxima at approximately
\SI{1.5}{\milli\second}, \SI{2.5}{\milli\second}, and
\SI{4.3}{\milli\second} are consistent with acoustic returns from
reflectors at distances of roughly \SI{25}{cm}, \SI{40}{cm}, and
\SI{70}{cm} from the array --- compatible with the dimensions of
the benchtop test setup, the housing of the array assembly, and
nearby fixtures.

The presence of stationary multipath has two implications. First,
the noise process at long range is correlated rather than white,
which affects the optimality of any matched-filter detection scheme
formulated under the white-noise assumption. The hardware engineer
who acquired the dataset reports that an initial naive attempt at
time-of-flight extraction --- bandpass filtering the
\texttt{Mic2} recordings at \SI{40}{kHz}, \SI{50}{kHz}, and
\SI{60}{kHz} separately and searching for a single envelope peak per
carrier --- was unsuccessful, and explicitly attributes the
difficulty to ``reflections from the environment''. The empirical
findings here are consistent with that statement: the reflections
are stationary, occupy the carrier bands, and remain visible in
target-absent recordings.
% NOTE: The engineer's statement is paraphrased from notes the
% author has shared in the conversation. If those notes become
% formally archived in the iCampus cloud or as a project
% technical report, replace this attribution with a \citep{} to
% that document. No external literature citation is required for
% this paragraph, since the relevant observation is documented
% within the project itself.
Second, the multipath component is deterministic in the sense that
it depends on the geometry of the test setup rather than on the
target, so any wavelet shrinkage that suppresses it must do so on
the basis of its time--frequency structure rather than its
randomness. The denoising scheme described in
Section~\ref{sec:wavelet-denoising} treats the multipath tail as
correlated additive content with energy concentrated in the same
carrier bands as the signal-of-interest, and relies on the time
localisation of the wavelet basis to separate the multipath energy
(spread across a long tail) from the burst energy of the direct
echo (concentrated in a short interval).

\subsection{The direct echo}
\label{sec:direct-echo}

\begin{figure}[tb]
    \centering
    \includegraphics[width=0.95\linewidth]{Pictures/fig_anatomy_rms.png}
    \caption{RMS amplitude in 200-sample (\SI{100}{\micro\second})
        windows for the \texttt{Box} at \SI{75}{cm} and $0\degree$
        (red) and for the target-absent reference (blue), both
        averaged over 100 trials. Logarithmic vertical axis. The
        two curves are essentially identical except at the direct
        echo near \SI{4.4}{\milli\second} and at the end-of-trial
        multi-bounce echo near \SI{8.7}{\milli\second}. The
        shaded grey region indicates the stationary-multipath
        interval; the dashed vertical line marks the blanking
        boundary at sample 2400.}
    \label{fig:anatomy-rms}
\end{figure}

Between samples \num{8800} and \num{9300}
(approximately \SI{4.4}{\milli\second}), Figure~\ref{fig:anatomy-rms}
shows a clear divergence between the measurement and reference RMS
profiles. The measurement RMS rises from approximately \num{4} to
above \num{30} \gls{adc} units; the reference remains flat near
the noise floor. This is the direct echo: the burst reflected from
the target back to the microphone. The arrival time of
\SI{4.4}{\milli\second} matches the expected geometric round trip
of $2 \cdot \SI{0.75}{\meter} / \SI{343}{\meter\per\second} =
\SI{4.37}{\milli\second}$ to within the resolution of the
\SI{100}{\micro\second} averaging window.

The spectral content of the direct echo, computed over samples
\num{8800} to \num{9300}, shows all three carriers present at
roughly equal magnitude, consistent with the simultaneous Tx
firing described in Section~\ref{sec:transmission-burst} and with
the acoustic propagation introducing negligible
frequency-dependent attenuation over the \SI{75}{cm} path. The
direct echo is, in this sense, a true ``copy'' of the transmitted
multi-frequency burst, scaled and time-shifted by the bistatic
geometry. This is the signal-of-interest for the matched filter
and the wavelet denoiser.

At the shortest protocol distance of \SI{25}{cm}, the direct echo
arrives at approximately \SI{1.5}{\milli\second}, immediately after
the transducer ringdown has ended at \SI{1}{\milli\second}. The
ringdown amplitude at \SI{0.3}{\milli\second}, near its peak, is
approximately \num{5} times the direct-echo peak at \SI{25}{cm};
by \SI{0.5}{\milli\second} the ringdown has dropped to below the
echo amplitude; and by \SI{0.8}{\milli\second} it is roughly one
order of magnitude smaller than the echo. The current blanking
length of \num{2400} samples (\SI{1.2}{\milli\second}), introduced
in the next section, removes both the burst and the residual
ringdown while leaving the entire \SI{25}{cm} echo intact.

\subsection{The end-of-trial peak}
\label{sec:end-peak}

A secondary feature in Figure~\ref{fig:anatomy-rms} is the small
RMS bump in the measurement at approximately
\SI{8.7}{\milli\second} (samples 17400 to 17800). The reference
shows no such feature. The position and the conditional appearance
of this peak make it worth investigating in its own right, both to
explain what it is and to determine whether the back end of the
trial should also be blanked.

\begin{figure}[tb]
    \centering
    \includegraphics[width=0.95\linewidth]{Pictures/fig_anatomy_endpeak.png}
    \caption{Investigation of the end-of-trial peak.
        \textbf{Top:} RMS profile in the interval
        \SIrange{7.0}{9.0}{\milli\second} for the \texttt{Box}
        target at all five protocol distances
        (\SIlist{25;50;75;100;125}{cm}), with the
        target-absent reference overlaid as a dashed black line.
        The large peak at \SI{7.3}{\milli\second} in the
        \SI{125}{cm} trace is the direct echo at long range. The
        smaller peaks near \SI{8.7}{\milli\second} in the
        \SI{50}{cm} and \SI{75}{cm} traces are the multi-bounce
        echo. \textbf{Bottom:} Spectrum of the direct echo
        (samples \numrange{8800}{9300}, red) and of the end peak
        (samples \numrange{17400}{17900}, blue) for the
        \texttt{Box} at \SI{75}{cm}, averaged over 100 trials.
        The direct echo contains all three carriers at
        approximately equal magnitude; the end peak is dominated
        by \SI{40}{kHz}, with the higher carriers progressively
        attenuated --- consistent with two acoustic
        reflection-plus-propagation cycles.}
    \label{fig:anatomy-endpeak}
\end{figure}

Figure~\ref{fig:anatomy-endpeak}
shows the end-of-trial region (samples 14000 to 18000) for the
\texttt{Box} target at all five protocol distances. At
\SI{125}{cm}, the entire end region is dominated by a large peak
at \SI{7.3}{\milli\second} --- this is the direct echo at long
range, where the geometric round-trip time of
$2 \cdot \SI{1.25}{\meter} / \SI{343}{\meter\per\second} =
\SI{7.29}{\milli\second}$ falls inside this window. At
\SI{100}{cm}, the direct echo arrives at \SI{5.83}{\milli\second}
(outside the end region) and the recording is at the noise floor.
At \SI{75}{cm} and \SI{50}{cm}, a smaller peak appears near
\SI{8.7}{\milli\second}, and the same peak is absent at
\SI{25}{cm} and \SI{100}{cm}.

The position of this secondary peak relative to the direct echo
matches a double round trip. For the \SI{75}{cm} case, the direct
echo is at \SI{4.37}{\milli\second} and the end peak is at
\SI{8.7}{\milli\second} --- a difference of
\SI{4.33}{\milli\second}, equal to a second round trip between
the array and the target. The mechanism is acoustic reciprocity:
the direct echo arrives at the receiver array, a fraction of it
re-radiates back toward the target, reflects, and returns to the
array a second time. The spectral content of the end peak,
visible in the lower panel of Figure~\ref{fig:anatomy-endpeak},
supports this interpretation: the \SI{40}{kHz} component
dominates, with \SI{50}{kHz} attenuated by approximately
\SI{8}{dB} and \SI{60}{kHz} attenuated by approximately
\SI{19}{dB} relative to the direct echo's spectrum at the same
frequencies. This monotonic frequency-dependent attenuation is
consistent with two acoustic reflection-plus-propagation cycles,
each of which absorbs higher frequencies preferentially.

The end peak is also object-dependent. At the same \SI{75}{cm}
distance and $0\degree$ angle, the end-peak RMS is approximately
\num{6.7} for \texttt{Box}, \num{4.5} for \texttt{Glas}, and only
\num{2.1} for \texttt{Dose} --- the latter being indistinguishable
from the noise floor. The cardboard \texttt{Box} reflects in a
sufficiently diffuse and broadband manner to support two
round-trip cycles; the metal \texttt{Dose} produces a strong first
echo but the geometry of subsequent reflections is unfavourable.

The implication for the blanking design is that the end-of-trial
region should \emph{not} be blanked. The multi-bounce signal is a
genuine acoustic event whose presence and shape carry information
about the target. It is potentially useful for classification, and
because the noise variance estimate is computed from the reference
recording (Section~\ref{sec:denoising-noise-estimation}) rather
than from the measurement trial, the end peak does not contaminate
the threshold used by the wavelet shrinkage step. No back-end
blanking is applied in this thesis.

\subsection{Implications for blanking and noise estimation}
\label{sec:anatomy-implications}

Three implications for the denoising pipeline follow from this
analysis.

First, the front-end blanking of \num{2400} samples
(\SI{1.2}{\milli\second}), described in
Section~\ref{sec:denoising-blanking}, removes the transmission
burst and the transducer ringdown without affecting any genuine
echo. The shortest protocol distance of \SI{25}{cm} corresponds to
a round-trip time of \SI{1.46}{\milli\second}, so any echo arrives
no earlier than sample \num{2920} --- comfortably after the
blanking boundary. By \SI{1.2}{\milli\second} the ringdown is
already an order of magnitude below the \SI{25}{cm} echo peak,
making the blanking redundant for echo preservation but valuable
for what it removes from any per-trial signal energy estimate.

Second, between the blanking boundary and approximately
\SI{5}{\milli\second}, the recording contains stationary multipath
from the test environment. The amplitude in this region is small
compared to the burst (typically \num{10} to \num{50} \gls{adc}
units) but well above the noise floor of \num{2} \gls{adc} units,
and it is correlated across trials. Any noise-variance estimate
that pools samples from this region in the time domain will be
biased upward by the multipath energy. The wavelet-domain
estimator used in Section~\ref{sec:denoising-noise-estimation}
sidesteps this bias by estimating the noise variance from a wavelet
detail level that contains no carrier-band content, isolating
broadband noise from the multipath energy concentrated near the
carriers.

Third, the noise process is not white: it is dominated, in the
operational time range of the recording, by correlated multipath
returns from the test setup. Matched-filter detection schemes
formulated under a white-noise assumption are therefore not
strictly optimal in this dataset, and the denoising stage acts in
part as a means of suppressing this correlated component before
the matched filter is applied. The wavelet basis is well suited to
this task because the multipath energy is spread across the
\SI{1}{\milli\second} to \SI{5}{\milli\second} interval --- a long
tail in the time domain, which corresponds to a small number of
large coefficients only at the coarsest wavelet scales --- whereas
the direct echo is a localised burst that produces large
coefficients at the scales matched to its envelope width.

%======================================================================
% Sources used in Section 3.2
%
% 1. Engineer's notes (text provided verbatim by the author in the
%    conversation):
%    - "Find time-of-flight ... not possible due to reflections
%      from the environment" — direct attribution of the multipath
%      problem to environmental reflections, used in
%      Section 3.2.3 (stationary multipath subsection).
%    - "The resonance frequency is around 39 kHz, so there should
%      not be any ringing" — engineer's a priori expectation,
%      referenced in Section 3.2.2 (ringdown subsection).
%
% 2. Knowles SPH18C3LM4H-1 datasheet \citep{Knowles2020}:
%    - Table 2: Fres = 39 kHz (used for the time constant
%      calculation in Section 3.2.2)
%    - Figure 7: Free Field Ultrasonic Response (used to estimate
%      Q ≈ 2.5 from the curve, with the uncertainty acknowledged
%      explicitly in the prose)
%
% 3. The dataset itself (verified by signal_anatomy.py):
%    - Region boundaries (0.3 ms, 1.0 ms, 5.0 ms)
%    - Spectral content of direct echo and end peak
%    - End-peak position and object dependence
%    - RMS comparison of measurement and reference
%
% TODO-AUTHOR items in Section 3.2:
%
% 1. If the engineer's notes are formally archived (in the iCampus
%    cloud or as a project technical report), replace the inline
%    attribution in Section 3.2.3 with a \citep{} to that document.
%    Currently the text refers to "the hardware engineer reports"
%    in prose.
%
% TODO-FIGURE items in Section 3.2:
%
% 1. fig_anatomy_full.png — generated by signal_anatomy.py
% 2. fig_anatomy_decay.png — generated by signal_anatomy.py
% 3. fig_anatomy_rms.png — generated by signal_anatomy.py
% 4. fig_anatomy_endpeak.png — generated by signal_anatomy.py
% (fig_anatomy_spectrogram.png is generated but not directly cited
%  in the section text; available as a supplementary figure or
%  for the wavelet-choice section later.)
%======================================================================
